% GENERATED FILE. Edit canonical lessons, then run npm run generate:books.
% canonical-source-hash: fnv1a64:0efbd9f8d04daf8e
% canonical-lessons: SA-C222-sadhanam, SA-C222-bharah, SA-C222-panyam, SA-C222-apanikah, SA-C222-patrika

\chapter{Instrument, Load, Goods for sale, Shopkeeper, Newspaper}
\label{ch:sa-instrument-load-goods-for-sale-shopkeeper-newspaper}
\hlchaptermodality{222}

\begin{chapteropening}
\textbf{By the end of this chapter:} \emph{I can say an instrument, a load, goods for sale, a shopkeeper and a newspaper.}

Complete the last lesson of chapter 222: I can say an instrument, a load, goods for sale, a shopkeeper and a newspaper.
\end{chapteropening}

\section[\texorpdfstring{\sk{साधनम्} (sādhanam)}{साधनम् (sādhanam)}]{\sk{साधनम्} (sādhanam) — an instrument, a means}

\label{lesson:SA-C222-sadhanam}



\begin{quote}
Before the new one: say the Sanskrit for a knot, then the Sanskrit for a small stool.
\end{quote}

\subsection*{You'll want to know: \sk{साधनम्}}
\textbf{\sk{साधनम्}} — \emph{sādhanam} — \textquotedblleft{}an instrument, a means\textquotedblright{}.

\begin{grammarlens}[title={What you've built}]
One more word for this chapter.
\end{grammarlens}

\subsection*{Your turn}
\begin{itemize}
\raggedright
  \item \emph{Say it:} \emph{sādhanam}
  \item \emph{Say it:} \emph{sādhanam}, once more
  \item \emph{Say it:} say \emph{sādhanam}
  \item \emph{Recall:} say the Sanskrit for a knot, then the Sanskrit for a small stool, then say \emph{sādhanam} again
\end{itemize}

\begin{tcolorbox}[breakable,colback=teal!4,colframe=teal!35!black,title={Before you move on}]
What does \sk{साधनम्} mean? (\textquotedblleft{}An instrument, a means\textquotedblright{}.) Say it once more.
\end{tcolorbox}
\section[\texorpdfstring{\sk{भारः} (bhāraḥ)}{भारः (bhāraḥ)}]{\sk{भारः} (bhāraḥ) — a load, a burden}

\label{lesson:SA-C222-bharah}



\begin{quote}
Before the new one: say the Sanskrit for a small stool, then the Sanskrit for an instrument.
\end{quote}

\subsection*{You'll want to know: \sk{भारः}}
\textbf{\sk{भारः}} — \emph{bhāraḥ} — \textquotedblleft{}a load, a burden\textquotedblright{}.

\begin{grammarlens}[title={What you've built}]
One more word for this chapter.
\end{grammarlens}

\subsection*{Your turn}
\begin{itemize}
\raggedright
  \item \emph{Say it:} \emph{bhāraḥ}
  \item \emph{Say it:} \emph{bhāraḥ}, once more
  \item \emph{Say it:} say \emph{bhāraḥ}
  \item \emph{Recall:} say the Sanskrit for a small stool, then the Sanskrit for an instrument, then say \emph{bhāraḥ} again
\end{itemize}

\begin{tcolorbox}[breakable,colback=teal!4,colframe=teal!35!black,title={Before you move on}]
What does \sk{भारः} mean? (\textquotedblleft{}A load, a burden\textquotedblright{}.) Say it once more.
\end{tcolorbox}
\section[\texorpdfstring{\sk{पण्यम्} (paṇyam)}{पण्यम् (paṇyam)}]{\sk{पण्यम्} (paṇyam) — goods for sale}

\label{lesson:SA-C222-panyam}



\begin{quote}
Before the new one: say the Sanskrit for an instrument, then the Sanskrit for a load.
\end{quote}

\subsection*{You'll want to know: \sk{पण्यम्}}
\textbf{\sk{पण्यम्}} — \emph{paṇyam} — \textquotedblleft{}goods for sale\textquotedblright{}.

\begin{grammarlens}[title={What you've built}]
One more word for this chapter.
\end{grammarlens}

\subsection*{Your turn}
\begin{itemize}
\raggedright
  \item \emph{Say it:} \emph{paṇyam}
  \item \emph{Say it:} \emph{paṇyam}, once more
  \item \emph{Say it:} say \emph{paṇyam}
  \item \emph{Recall:} say the Sanskrit for an instrument, then the Sanskrit for a load, then say \emph{paṇyam} again
\end{itemize}

\begin{tcolorbox}[breakable,colback=teal!4,colframe=teal!35!black,title={Before you move on}]
What does \sk{पण्यम्} mean? (\textquotedblleft{}Goods for sale\textquotedblright{}.) Say it once more.
\end{tcolorbox}
\section[\texorpdfstring{\sk{आपणिकः} (āpaṇikaḥ)}{आपणिकः (āpaṇikaḥ)}]{\sk{आपणिकः} (āpaṇikaḥ) — a shopkeeper}

\label{lesson:SA-C222-apanikah}



\begin{quote}
Before the new one: say the Sanskrit for a load, then the Sanskrit for goods for sale.
\end{quote}

\subsection*{You'll want to know: \sk{आपणिकः}}
\textbf{\sk{आपणिकः}} — \emph{āpaṇikaḥ} — \textquotedblleft{}a shopkeeper\textquotedblright{}.

\begin{grammarlens}[title={What you've built}]
One more word for this chapter.
\end{grammarlens}

\subsection*{Your turn}
\begin{itemize}
\raggedright
  \item \emph{Say it:} \emph{āpaṇikaḥ}
  \item \emph{Say it:} \emph{āpaṇikaḥ}, once more
  \item \emph{Say it:} say \emph{āpaṇikaḥ}
  \item \emph{Recall:} say the Sanskrit for a load, then the Sanskrit for goods for sale, then say \emph{āpaṇikaḥ} again
\end{itemize}

\begin{tcolorbox}[breakable,colback=teal!4,colframe=teal!35!black,title={Before you move on}]
What does \sk{आपणिकः} mean? (\textquotedblleft{}A shopkeeper\textquotedblright{}.) Say it once more.
\end{tcolorbox}
\section[\texorpdfstring{\sk{पत्रिका} (patrikā)}{पत्रिका (patrikā)}]{\sk{पत्रिका} (patrikā) — a newspaper, a magazine}

\label{lesson:SA-C222-patrika}



\begin{quote}
Before the new one: say the Sanskrit for goods for sale, then the Sanskrit for a shopkeeper.
\end{quote}

\subsection*{You'll want to know: \sk{पत्रिका}}
\textbf{\sk{पत्रिका}} — \emph{patrikā} — \textquotedblleft{}a newspaper, a magazine\textquotedblright{}.

\begin{grammarlens}[title={What you've built}]
One more word for this chapter.
\end{grammarlens}

\subsection*{Your turn}
\begin{itemize}
\raggedright
  \item \emph{Say it:} \emph{patrikā}
  \item \emph{Say it:} \emph{patrikā}, once more
  \item \emph{Say it:} say \emph{patrikā}
  \item \emph{Recall:} say the Sanskrit for goods for sale, then the Sanskrit for a shopkeeper, then say \emph{patrikā} again
\end{itemize}

\begin{tcolorbox}[breakable,colback=teal!4,colframe=teal!35!black,title={Before you move on}]
What does \sk{पत्रिका} mean? (\textquotedblleft{}A newspaper, a magazine\textquotedblright{}.) Say it once more.
\end{tcolorbox}
